% mit großem Interesse habe ich Ihre Stellenausschreibung für die Position als Junior Softwareentwickler im Bereich Outputmanagement gelesen. Als angehender Absolvent der Angewandten Mathematik und Informatik (AMI) sowie ausgebildeter Mathematisch-technischer Softwareentwickler habe ich fundierte Kenntnisse in der Softwareentwicklung, insbesondere in C\# und der objektorientierten Programmierung, sowie praktische Erfahrungen in der Analyse und Optimierung von Softwarelösungen gesammelt. Diese Kombination aus theoretischem Wissen und praktischen Erfahrungen macht mich zu einem idealen Kandidaten für diese Position.

% mit großem Interesse habe ich Ihre Stellenausschreibung für die Position als Junior Software Entwickler im Bereich Outputmanagement gelesen. Als angehender Absolvent der Angewandten Mathematik und Informatik sowie ausgebildeter Mathematisch-technischer Softwareentwickler habe ich fundierte Kenntnisse in der Softwareentwicklung, insbesondere in C\# und der objektorientierten Programmierung sowie praktische Erfahrungen mit Datenbankmanagementsystemen wie MySQL und MariaDB gesammelt.

% Mit großem Interesse habe ich Ihre Stellenausschreibung für die Position als Junior Softwareentwickler im Bereich Outputmanagement gelesen. Als angehender Absolvent der Angewandten Mathematik (AMI) und Informatik sowie ausgebildeter Mathematisch-technischer Softwareentwickler bringe ich fundierte Kenntnisse in der Softwareentwicklung, insbesondere in der objektorientierten Programmierung, mit. Zudem habe ich praktische Erfahrungen mit Datenbankmanagementsystemen wie MySQL und MariaDB gesammelt.

% Mit großer Begeisterung habe ich Ihre Stellenausschreibung für die Position des Trainee Product Managers gelesen. Als angehender Absolvent der Angewandten Mathematik und Informatik sowie ausgebildeter Mathematisch-technischer Softwareentwickler bringe ich fundiertes technisches Verständnis, analytische Fähigkeiten und eine hohe IT-Affinität mit. Die Aussicht, in einem dynamischen Umfeld umfassende Einblicke in die Produktentwicklung zu erhalten und dabei Optimierungsprozesse aktiv mitzugestalten, begeistert mich besonders. Ich freue mich darauf, mein Wissen in spannenden Projekten einzubringen und und gemeinsam mit Ihrem Team an der erfolgreichen Umsetzung der Produktstrategien zu arbeiten.

% Mit großer Begeisterung habe ich Ihre Stellenausschreibung für die Position des Traineeprogramms in der Softwareentwicklung gelesen. Als angehender Absolvent der Angewandten Mathematik und Informatik (AMI) sowie ausgebildeter Mathematisch-technischer Softwareentwickler (MATSE) bringe ich fundierte Kenntnisse in der objektorientierten Programmierung mit, insbesondere in Java, und bin motiviert, mein Wissen in einem dynamischen Umfeld weiter zu vertiefen.

% mit großer Begeisterung habe ich Ihre Stellenausschreibung für die Position des Trainee Softwareentwicklung gelesen. Besonders reizt mich die Möglichkeit, in Ihrem Unternehmen meine Kenntnisse in der Softwareentwicklung in der Praxis anzuwenden, zu erweitern und mich mit neuen Technologien wie Microsoft Dynamics 365 Business Central auseinanderzusetzen. hr umfassendes Trainee-Programm, das zielgerichtete Schulungen und Trainings bietet, sowie die Arbeit an individuellen Softwarelösungen, spricht mich sehr an.

% Mit großer Begeisterung habe ich Ihre Stellenausschreibung für die Position des Trainee Softwareentwicklung gelesen. Besonders schätze ich die Möglichkeit, meine Kenntnisse in der Softwareentwicklung in Ihrem Unternehmen praktisch einzubringen und weiterzuentwickeln. Ich freue mich darauf, mich intensiv mit neuen Technologien wie Microsoft Dynamics 365 Business Central und der AL-Programmierung auseinanderzusetzen sowie moderne Projektmethoden zu erlernen und in der Praxis umzusetzen.


% mit großer Begeisterung habe ich Ihre Stellenausschreibung für die Position des Trainee Softwareentwicklung und Testautomatisierung gelesen. Besonders schätze ich die Möglichkeit, in einem intensiven Bootcamp die Methoden der modernen Software-Qualitätssicherung zu erlernen und direkt in die Praxis umzusetzen.


% mit großer Begeisterung habe ich Ihre Stellenausschreibung für die Position des Trainee Softwareentwicklung gelesen. Besonders reizt mich die Möglichkeit, in Ihrem Unternehmen meine Kenntnisse in der Softwareentwicklung praktisch anzuwenden und weiterzuentwickeln. Ich freue mich darauf, mich intensiv mit neuen Technologien wie Microsoft Dynamics 365 Business Central und der AL-Programmierung auseinanderzusetzen sowie moderne Projektmethoden kennenzulernen. Ihr umfassendes Trainee-Programm, das zielgerichtete Schulungen und Trainings bietet, sowie die Arbeit an individuellen Softwarelösungen spricht mich sehr an. In einem dynamischen Umfeld zu lernen und aktiv zur Digitalisierung im Mittelstand beizutragen, motiviert mich besonders.


% mit großem Interesse habe ich Ihre Stellenanzeige für die Position des Junior Java Entwicklers gelesen. Die Möglichkeit, aktiv an der digitalen Transformation der Versicherungsbranche mitzugestalten und dabei moderne IT-Technologien zu realisieren, reizt mich sehr. Die Aussicht, in einem hochqualifizierten und aufgeschlossenen Projektteam an der Entwicklung nachhaltiger Software-Lösungen zu arbeiten, entspricht genau meinen beruflichen Zielen und Interessen.

% mit großem Interesse habe ich Ihre Stellenausschreibung für die Position des Junior Java Entwicklers gelesen. Die Möglichkeit, aktiv am digitalen Wandel der Versicherungsbranche mitzuwirken und innovative Software-Lösungen zu entwickeln, begeistert mich sehr.

% mit großem Interesse habe ich Ihre Stellenausschreibung für die Position des Junior Softwareentwicklers gelesen. Die Möglichkeit, aktiv am Wandel der digitalen Technologien mitzuwirken und dabei innovative Software-Lösungen zu entwickeln, begeistert mich sehr.

% mit großem Interesse habe ich Ihre Stellenausschreibung für die Position des Junior Softwareentwickler gelesen. Die Aussicht, innovative Softwarelösungen mit Microsoft-Technologien wie C\# und modernen agilen Methoden zu entwickeln, begeistert mich sehr.

% mit großer Begeisterung habe ich erfahren, dass Sie auf der Suche nach einem Junior Softwareentwickler sind. Die Möglichkeit, Softwarelösungen zu entwickeln, die die Geschäftsprozesse Ihrer Kunden verbessern, und dabei in einem dynamischen, internationalen Umfeld zu arbeiten, schätze ich sehr. Ich freue mich darauf, meine Fähigkeiten in einem Team einzubringen, das den gesamten Entwicklungsprozess – von der ersten Anforderung bis zur endgültigen Auslieferung – mitgestaltet und dabei stets die höchste Qualität anstrebt.

% mit großer Begeisterung habe ich erfahren, dass Sie auf der Suche nach einem Junior Softwareentwickler sind. Die Aussicht, meine Kenntnisse in der Entwicklung innovativer Softwarelösungen einzubringen und dabei eng mit einem internationalen Team zusammenzuarbeiten, begeistert mich sehr. Besonders schätze ich die Möglichkeit, den gesamten Entwicklungsprozess – von der Anforderungsanalyse bis zur finalen Auslieferung – aktiv mitzugestalten.

% mit großem Interesse habe ich Ihre Stellenausschreibung für die Position des Junior Software Developer .NET gelesen. Die Möglichkeit, digitale und webbasierte Lösungen für internationale Kunden zu entwickeln und dabei moderne Microsoft-Technologien wie C\# in agilen Umgebungen zu nutzen, spricht mich sehr an.

% mit großem Interesse habe ich Ihre Stellenausschreibung für die Position des Junior IT Anwendungsentwicklers gelesen. meine Kenntnisse in der Anwendungsentwicklung zu vertiefen und gleichzeitig an der Wartung und Optimierung bestehender Anwendungen mitzuwirken, spricht mich sehr an. 

% mit großem Interesse habe ich Ihre Ausschreibung für die Trainee-Stelle in der Softwareentwicklung bei Buhl gelesen. Die Aussicht, in einem innovativen Umfeld Softwarelösungen zu entwickeln und gleichzeitig meine Kenntnisse in C\#, JavaScript und PHP weiter auszubauen, begeistert mich sehr.

% mit großem Interesse habe ich Ihre Ausschreibungen für die Positionen als Junior Softwareentwickler, Junior C\# Entwickler und Junior Java Entwickler gelesen. Die Aussicht, in einem innovativen Team Softwarelösungen zu entwickeln und meine Kenntnisse in der Softwareentwicklung weiter auszubauen, begeistert mich sehr.

% mit großem Interesse habe ich Ihre Ausschreibung für die Position als Junior Softwaretester gelesen. Die Aussicht, in einem innovativen Scrum-Team Softwarelösungen mitzuwirken und meine Kenntnisse in der Softwareentwicklung weiter auszubauen, begeistert mich sehr.

% mit großem Interesse habe ich Ihre Ausschreibung für die Trainee-Position in der Softwareentwicklung gelesen. Besonders schätze ich die Möglichkeit, alle Phasen der Anwendungsentwicklung in einem dynamischen Team kennenzulernen und eigenverantwortlich an Projekten mitzuwirken. 

% mit großer Begeisterung habe ich die Stellenausschreibung für das Junior Expert Programm IT gelesen, da es die perfekte Gelegenheit bietet, meine Leidenschaft für agile Softwareentwicklung, Cloud-Technologien und innovative IT-Lösungen in einem interdisziplinären Umfeld zu vertiefen. Als Absolvent der Angewandten Mathematik und Informatik (AMI) sowie ausgebildeter Mathematisch-technischer Softwareentwickler (MATSE) bringe ich fundierte Kenntnisse in der Softwareentwicklung sowie in Web-Technologien  mit, die ich gerne in Ihrem Team weiterentwickeln möchte.
% mit großem Interesse habe ich Ihre Stellenausschreibung gelesen und bin überzeugt, dass die Position als Junior Software Entwickler eine perfekte Gelegenheit bietet, meine Kenntnisse in Java und Webtechnologien weiter auszubauen und aktiv an innovativen Softwareprojekten mitzuwirken. Nachdem ich mein Masterstudium in Mathematik aufgrund finanzieller Herausforderungen nicht abschließen konnte, entschied ich mich bewusst für ein duales Studium in Angewandter Mathematik und Informatik (AMI), das mir half, sowohl meine praktischen als auch theoretischen Kenntnisse in der Softwareentwicklung zu vertiefen.

% mit großem Interesse habe ich Ihre Stellenausschreibung gelesen und bin überzeugt, dass die Position als Junior Software- Testen Entwickler eine perfekte Gelegenheit bietet, meine Kenntnisse in Java und Webtechnologien weiter auszubauen aktiv an der Konzeption und Entwicklung innovativer Web- und Intranet-Anwendungen für Ihre Kunden mitzuwirken.
% {\fontsize{10.4pt}{12.8pt}\selectfont  Nachdem ich mein Masterstudium in Mathematik aufgrund finanzieller Herausforderungen nicht abschließen konnte, entschied ich mich bewusst für ein duales Studium in Angewandter Mathematik und Informatik (AMI), das mir half, sowohl meine theoretischen als auch praktischen Kenntnisse in der Softwareentwicklung zu vertiefen.}

% mit großem Interesse habe ich Ihre Stellenausschreibung gelesen und bin überzeugt, dass die Position als Junior Softwareentwickler / Testentwickler eine perfekte Gelegenheit bietet, meine Kenntnisse in Java weiter auszubauen und aktiv zur Qualitätssicherung komplexer Softwarelösungen im öffentlichen Sektor beizutragen.Nachdem ich mein Masterstudium in Mathematik aufgrund finanzieller Herausforderungen nicht abschließen konnte, entschied ich mich bewusst für ein duales Studium in Angewandter Mathematik und Informatik (AMI), das mir half, sowohl meine theoretischen als auch praktischen Kenntnisse in der Softwareentwicklung zu vertiefen.

mit großem Interesse habe ich Ihre Stellenausschreibung gelesen und bin überzeugt, dass die Position als Analyst für Energie-Stammdaten eine perfekte Gelegenheit bietet, meine Kenntnisse in der Analyse und Auswertung großer Datenmengen optimal einzubringen und aktiv zur weiteren Gestaltung der Energiewende beizutragen.

Nachdem ich mein Masterstudium in Mathematik aufgrund finanzieller Herausforderungen nicht abschließen konnte, entschied ich mich bewusst für ein duales Studium in Angewandter Mathematik und Informatik (AMI), das mir half, sowohl meine theoretischen als auch praktischen Kenntnisse in der Softwareentwicklung zu vertiefen. Während meines erfolgreich abgeschlossenen AMI-Studiums habe ich tiefgehende Erfahrungen im Umgang mit SQL-basierten \href{https://github.com/montahaee/DB/tree/master/ex_am_cise/docker}{\textcolor{green!37!black}{Datenbanken}} (MySQL, MariaDB) und im \href{https://github.com/montahaee/DS}{\textcolor{green!37!black}{Data Science}} bzw. in der \href{https://montahaee.github.io/assets/pdf/en/projects/heart_disease_Prediction.pdf}{\textcolor{green!37!black}{Datenanalyse}} mit Hilfe Python-Biblothken wie pandas, sklearn und TensorFlow gesammelt, was mich ideal für die von Ihnen genannten Anforderungen qualifiziert. Insbesondere die Datenbearbeitung und -auswertung mit der gennanten Bilotheken ermöglichte mir eine effiziente Handhabung großer Datenmengen.

Darüber hinaus habe ich in meiner \href{https://montahaee.github.io/assets/pdf/de/projects/JLL1_final.pdf}{\textcolor{green!37!black}{Seminararbeit}} mathematisch anlysiert, wie die Ressourcenbelastung bei der Verarbeitung hoher Datenvolumina durch Dimensionalitätsreduktionstechniken (u. a. Johnson-Lindenstrauss-Lemma) in Vergleich zu der Hauptkomponentenanalyse (PCA) minimiert werden kann. Diese Arbeit zielte darauf ab, die Datenanalyse effizienter zu gestalten, ohne signifikanten Informationsverlust. 

Auch in meiner \href{https://montahaee.github.io/assets/pdf/de/projects/myFH_Bachelorthesis.pdf}{\textcolor{green!37!black}{Bachelorarbeit}} im Bereich Machine Learning konnte ich meine Kenntnisse vertiefen: Ich untersuchte eine datengetriebene exakte Physik für wissenschaftlisch machinelles Lernrn , wodurch ich ein umfassendes Verständnis für datengetriebene Analyseprozesse sowohl mathematisch als auch stochastisch erlangte und somti meine Fähigkeit, komplexe Datenzusammenhänge präzise zu modellieren und auszuwerten, verstärkte.
%  Ihr Team beizusteuern und gemeinsam zur Weiterentwicklung  und innovative Lösungen der Energiedatenanalyse beizutragen. Über die Möglichkeit, meine Fähigkeiten und Motivation in einem persönlichen Gespräch näher vorzustellen, würde ich mich sehr freuen.

% Weiterhin habe ich mein Wissen im objektiven Datenmanagement vor allem durch meine Mitarbeit an der Integration von Forschungsdatenmodellen auf der \href{https://coscine.rwth-aachen.de/login?redirect=/}{\textcolor{green!37!black}{Coscine}}-Plattform erweitert, was mir einen tiefen Einblick in agile Arbeitsmethoden und moderne Cloud-Technologien ermöglichte.

Weiterhin konnte ich durch meine Mitarbeit an der Integration von Forschungsdatenmodellen auf der \href{https://coscine.rwth-aachen.de/login?redirect=/}{\textcolor{green!37!black}{Coscine}}-Plattform nicht nur wertvolle Einblicke in Objektdatenbanken und Cloud-Technologien, sondern auch in die Zusammenarbeit in interdisziplinären Teams gewinnen. Die Arbeit in einem agilen Umfeld, das iterative und inkrementelle Prozesse kombiniert, ermöglichte es mir, flexibel auf Veränderungen zu reagieren und die Software qualitativ zu verbessern.
% Nachdem ich mein Masterstudium in Mathematik aufgrund finanzieller Herausforderungen nicht abschließen konnte, entschied ich mich bewusst für ein duales Studium in Angewandter Mathematik und Informatik (AMI), das mir half, sowohl meine theoretischen als auch praktischen Kenntnisse in der Softwareentwicklung zu vertiefen. In meinem erfolgreich abgeschlossenen AMI-Studium konnte ich als anstrebender Mathematisch-technischer Softwareentwickler im Rahmen verschiedener Projekte und Module wertvolle \href{https://montahaee.github.io/repositories/}{\textcolor{cyan!50!black}{Erfahrungen}} in der objektorientierten Programmierung, der Optimierung von Softwarekomponenten sowie in der Entwicklung von Java-, C\#- und Python-basierten Applikationen sammeln. Ein besonderer Fokus lag dabei auf der Anwendung von \href{https://de.wikipedia.org/wiki/Entwurfsmuster}{\textcolor{cyan!50!black}{Entwurfsmustern}}, die zur Strukturierung, Flexibilität und Wiederverwendbarkeit des Codes beitragen, sowie auf der Durchführung von Unit-Tests, um die Funktionalität einzelner Softwarekomponenten zu sichern und robuste, wartbare Softwarelösungen zu gewährleisten. Die Entwicklung automatisierter Systemintegrationstests für sicherheitskritische Anwendungen sehe ich daher als spannende Aufgabe, die meinem Interesse an strukturierten Testansätzen und meinen methodischen Fähigkeiten sehr entgegenkommt.

% Besonders intensiv vertiefte ich meine Kenntnisse in Java während der \href{https://github.com/montahaee/Informatics}{\textcolor{green!37!black}{Vorbereitung}} und praktischen Durchführung der IHK-Abschlussprüfung zur Entwicklung eines \href{https://github.com/montahaee/SS2023/tree/main/softwareProduct}{\textcolor{green!37!black}{Softwareprodukts}}. Dabei wurde das \textit{Producer-Consumer-Dataflow }Muster in einem multithreaded Kontext eingesetzt, um Datenflüsse parallelisiert zu verarbeiten und eine optimale Ressourcennutzung zu gewährleisten.


% mit großem Interesse habe ich von der Möglichkeit erfahren, Ihre Kunden durch die Entwicklung moderner Web- und Intranetanwendungen zu unterstützen. Die Aussicht, in Projekten die Struktur und Architektur von Software zu gestalten und meine Java-Kenntnisse dabei gezielt zu vertiefen, motiviert mich sehr. Nachdem ich mein Masterstudium in Mathematik aufgrund finanzieller Herausforderungen nicht abschließen konnte, entschied ich mich bewusst für ein duales Studium in Angewandter Mathematik und Informatik (AMI), das mir half, sowohl meine praktischen als auch theoretischen Kenntnisse in der Softwareentwicklung zu vertiefen.

% In meinem erfolgreich abgeschlossenen AMI-Studium konnte ich als anstrebender Mathematisch-technischer Softwareentwickler im Rahmen verschiedener Projekte und Module wertvolle \href{https://montahaee.github.io/repositories/}{\textcolor{cyan!50!black}{Erfahrung}} in der objektorientierten Programmierung, der Optimierung von Softwarekomponenten, sowie in der Entwicklung von Java-, C\#- und Python-basierten Applikationen erworben. Ein besonderer Fokus lag dabei auf der Anwendung von \href{https://montahaee.github.io/blog/2023/dive-into-design-pattern/}{\textcolor{cyan!50!black}{Entwurfsmustern}}, die zur Strukturierung, Flexibilität und Wiederverwendbarkeit des Codes beitragen, sowie auf der Durchführung von Unit-Tests, um die Funktionalität einzelner Softwarekomponenten zu sichern und robuste, wartbare Softwarelösungen zu gewährleisten. Darüber hinaus habe ich mein Wissen im objektiven Datenmanagement vor allem durch meine Mitarbeit an der Integration von Forschungsdatenmodellen auf der \href{https://coscine.rwth-aachen.de/login?redirect=/}{\textcolor{green!37!black}{Coscine}}-Plattform erweitert, was mir einen tiefen Einblick in agile Arbeitsmethoden und moderne Cloud-Technologien ermöglichte.

% In meinem erfolgreich abgeschlossenen AMI-Studium konnte ich als anstrebender Mathematisch-technischer Softwareentwickler im Rahmen verschiedener Projekte und Module wertvolle \href{https://montahaee.github.io/repositories/}{\textcolor{cyan!50!black}{Erfahrung}} in der objektorientierten Programmierung, der Optimierung von Softwarekomponenten, sowie in der Entwicklung von Java-, C\#- und Python-basierten Applikationen erworben. Ein besonderer Fokus lag dabei auf der Anwendung von \href{https://montahaee.github.io/blog/2023/dive-into-design-pattern/}{\textcolor{cyan!50!black}{Entwurfsmustern}}, die zur Strukturierung, Flexibilität und Wiederverwendbarkeit des Codes beitragen, sowie auf der Durchführung von Unit-Tests, um die Funktionalität einzelner Softwarekomponenten zu sichern und robuste, wartbare Softwarelösungen zu gewährleisten. 
% Meine Vertifung in Java fand im Zug der \href{https://github.com/montahaee/Informatics}{\textcolor{violet}{Vorbereitung}} und praktischen Durchführung der IHK-Abschlussprüfung zur Entwicklung eines \href{https://github.com/montahaee/SS2023/tree/main/softwareProduct}{\textcolor{violet}{Softwareprodukts}} habe ich intensiv mit Java gearbeitet und das \textit{Producer-Consumer-Dataflow-Pattern} in einem multithreaded Kontext eingesetzt, um Datenflüsse parallelisiert zu verarbeiten und eine optimale Ressourcennutzung zu gewährleisten. 

% In meinem erfolgreich abgeschlossenen AMI-Studium konnte ich als anstrebender Mathematisch-technischer Softwareentwickler im Rahmen verschiedener Projekte und Module wertvolle \href{https://montahaee.github.io/repositories/}{\textcolor{cyan!50!black}{Erfahrung}} in der objektorientierten Programmierung, der Optimierung von Softwarekomponenten sowie in der Entwicklung von Java-, C\#- und Python-basierten Applikationen sammeln. Ein besonderer Fokus lag dabei auf der \href{https://montahaee.github.io/blog/2023/dive-into-design-pattern/}{\textcolor{cyan!50!black}{Anwendung von Entwurfsmustern}}, die zur Strukturierung, Flexibilität und Wiederverwendbarkeit des Codes beitragen, sowie auf der Durchführung von Unit-Tests, um die Funktionalität einzelner Softwarekomponenten zu sichern und robuste, wartbare Softwarelösungen zu gewährleisten.\\

% Eine intensive Vertiefung in Java erfolgte im Zuge der \href{https://github.com/montahaee/Informatics}{\textcolor{violet}{Vorbereitung}} und praktischen Durchführung der IHK-Abschlussprüfung zur Entwicklung eines \href{https://github.com/montahaee/SS2023/tree/main/softwareProduct}{\textcolor{violet}{Softwareprodukts}}. In diesem Projekt habe ich das \textit{Producer-Consumer-Dataflow-Pattern} in einem multithreaded Kontext eingesetzt, um Datenflüsse parallelisiert zu verarbeiten und eine optimale Ressourcennutzung zu gewährleisten.


% mit großem Interesse habe ich Ihre Ausschreibung für die Position als Fachinformatiker gelesen. Die Aussicht, in einem innovativen Scrum-Team Softwarelösungen zu entwickeln und meine Kenntnisse in der Softwareentwicklung weiter auszubauen, begeistert mich sehr.

% mit großer Begeisterung habe ich Ihre Ausschreibung für die Trainee-Position in der Softwareentwicklung bei HR WORKS gelesen. Die Möglichkeit, in einem innovativen und wachstumsstarken HR-Tech-Unternehmen mitzuwirken und die Entwicklung moderner Softwarelösungen im Bereich der digitalen Reisekostenabrechnung voranzutreiben, spricht mich sehr an.

% Während meines Studiums der Angewandten Mathematik und Informatik sowie meiner Ausbildung zum Mathematisch-technischen Softwareentwickler konnte ich wertvolle Erfahrungen in der Softwareentwicklung sammeln. Insbesondere habe ich Kenntnisse in der Programmierung mit Java, C\# und SQL erworben und in verschiedenen Projekten vertieft. Dazu gehören unter anderem die Entwicklung eines Gitter-Generators unter Verwendung des Adapter-Patterns sowie die Gestaltung eines UML-konformen Workflow-Managements für additive Fertigungsprozesse, inspiriert durch das Mediator-Pattern. Diese Projekte haben mir gezeigt, wie wichtig eine durchdachte Softwarearchitektur und der Einsatz geeigneter Design Patterns für eine saubere und wartbare Codebasis sind. 

% Des Weiteren habe ich praktische Erfahrungen in der Nutzung und Integration von Cloud-Services durch die Mitwirkung an der Integration von Forschungsdatenmodellen auf der Coscine-Plattform gesammelt. Diese Erfahrung hat mir gezeigt, wie wichtig es ist, Cloud-basierte Technologien effektiv einzusetzen, um Datenmanagement und wissenschaftliche Forschung zu unterstützen. 

% Während meines Studiums der Angewandten Mathematik und Informatik sowie meiner Ausbildung zum Mathematisch-technischen Softwareentwickler konnte ich umfassende Erfahrungen in der Softwareentwicklung sammeln. Besonders ausgeprägt sind meine Kenntnisse in der Programmierung mit Java, C\# und Python, die ich durch verschiedene Projekte vertieft habe. Dazu zählen die Entwicklung eines \href{https://montahaee.github.io/de/projects/1_project/}{\textcolor{violet}{Gitter-Generators}} unter Einsatz des Adapter-Patterns und die \href{https://montahaee.github.io/de/projects/3_project/}{\textcolor{violet}{Entwurf}} eines UML-konformen Workflow-Managements für additive Fertigungsprozesse, inspiriert durch das Mediator-Pattern. Diese Erfahrungen haben mir die Bedeutung einer durchdachten Softwarearchitektur und des gezielten Einsatzes von Design Patterns für eine saubere und wartbare Codebasis vor Augen geführt.

% Des Weiteren habe ich praktische Erfahrungen in der Nutzung und Integration von Cloud-Services durch die Mitwirkung an der Integration von Forschungsdatenmodellen auf der \href{https://coscine.rwth-aachen.de/login?redirect=/}{\textcolor{violet}{Coscine}}-Plattform gesammelt. Diese Aufgabe hat mir verdeutlicht, wie wichtig es ist, Cloud-basierte Technologien effektiv einzusetzen, um Datenmanagement und wissenschaftliche Forschung zu unterstützen. 

% Zusätzlich habe ich mich intensiv mit Cloud-Services auseinandergesetzt, insbesondere durch meine Mitwirkung an der Integration von Forschungsdatenmodellen auf der Coscine-Plattform. Diese Aufgabe hat mir verdeutlicht, wie essenziell der effektive Einsatz von Cloud-Technologien für das Datenmanagement und die Unterstützung wissenschaftlicher Forschung ist.

% Während meines AMI-Studiums und meiner MATSE-Ausbildung habe ich umfassende Kenntnisse in der Java-Programmierung erworben, insbesondere im Rahmen der \href{https://github.com/montahaee/Informatics}{\textcolor{violet}{Vorbereitung}} und praktischen Durchführung der IHK-Abschlussprüfung zur Entwicklung eines \href{https://github.com/montahaee/SS2023/tree/main/softwareProduct}{\textcolor{violet}{Softwareprodukts}}. Dabei habe ich intensiv mit Java gearbeitet und praktische Anwendungen entwickelt, unter anderem durch den Einsatz des \textit{Producer-Consumer- Dataflow-Patterns} in einem multithreaded Kontext, um Datenflüsse parallelisiert zu verarbeiten und gleichzeitig eine optimale Ressourcennutzung zu gewährleisten. 



% Zusätzlich habe ich praktische Erfahrungen in der Implementierung und Optimierung von Softwarekomponenten gesammelt, wie etwa in der Entwicklung eines \href{https://montahaee.github.io/de/projects/1_project/}{\textcolor{violet}{Gitter-Generators}} und dem \href{https://montahaee.github.io/de/projects/3_project/}{\textcolor{violet}{Entwurf}} eines UML-konformen Workflow-Managements für additive Fertigungsprozesse. In diesen Projekten wurden das Adapter bzw. das Mediator Pattern angewandt. Diese Projekte haben mir verdeutlicht, wie wichtig Clean Code und eine durchdachte Softwarearchitektur sind, und sie haben meine Leidenschaft für die Softwareentwicklung weiter gestärkt.

% Während meines Studiums der Angewandten Mathematik und Informatik sowie meiner Ausbildung zum Mathematisch-technischen Softwareentwickler habe ich praktische Erfahrungen in verschiedenen Programmiersprachen gesammelt. So habe ich beispielsweise die Integration von Forschungsdatenmodellen in Python auf der \href{https://coscine.rwth-aachen.de/login?redirect=/}{\textcolor{green!37!black}{Coscine}}-Plattform unterstützt. Bei der Entwicklung eines \href{https://montahaee.github.io/de/projects/1_project/}{\textcolor{green!37!black}{Gitter-Generators}} und eines \href{https://github.com/montahaee/SS2023/tree/main/softwareProduct}{\textcolor{green!37!black}{Softwareprodukts}} im Rahmen meiner IHK-Abschlussprüfung habe ich sowohl Black-Box-Tests als auch White-Box-Tests erstellt, um die Funktionalität und Qualität der Software sicherzustellen.


% Während meines Studiums der Angewandten Mathematik und Informatik sowie meiner Ausbildung zum Mathematisch-technischen Softwareentwickler habe ich umfassende Kenntnisse in der Java-Programmierung erworben, insbesondere im Rahmen der \href{https://github.com/montahaee/Informatics}{\textcolor{violet}{Vorbereitung}} auf praktischen Durchführung der IHK-Abschlussprüfung zur Entwicklung eines \href{https://github.com/montahaee/SS2023/tree/main/softwareProduct}{\textcolor{violet}{Softwareprodukts}}. Dabei habe ich intensiv mit Java gearbeitet und praktische Anwendungen entwickelt, unter anderem durch den Einsatz des \textit{Producer-Consumer- Dataflow-Patterns} in einem multithreaded Kontext, um Datenflüsse parallelisiert zu verarbeiten und gleichzeitig eine optimale Ressourcennutzung zu gewährleisten. Zusätzlich konnte ich durch die Unterstützung bei der Integration von Forschungsdatenmodellen auf der \href{https://coscine.rwth-aachen.de/login?redirect=/}{\textcolor{green!37!black}{Coscine}}-Plattform wertvolle Erfahrungen im Bereich Datenmanagement und Cloud-Technologien sammeln.


% Während meines Studiums der Angewandten Mathematik und Informatik sowie meiner Ausbildung zum Mathematisch-technischen Softwareentwickler habe ich praktische Erfahrung in der Implementierung und Optimierung von Softwarekomponenten gesammelt, wie etwa in der Entwicklung eines \href{https://montahaee.github.io/de/projects/1_project/}{\textcolor{violet}{Gitter-Generators}} und dem \href{https://montahaee.github.io/de/projects/3_project/}{\textcolor{violet}{Entwurf}} eines UML-konformen Workflow-Managements für additive Fertigungsprozesse. In diesen Projekten wurden das Adapter bzw. das Mediator Pattern angewandt.

% Darüber hinaus habe ich wertvolle Kenntnisse in der Java-Programmierung erworben, insbesondere im Rahmen der \href{https://github.com/montahaee/Informatics}{\textcolor{violet}{Vorbereitung}} auf praktischen Durchführung der IHK-Abschlussprüfung zur Entwicklung eines \href{https://github.com/montahaee/SS2023/tree/main/softwareProduct}{\textcolor{violet}{Softwareprodukts}}. Dabei habe ich intensiv mit Java gearbeitet und praktische Anwendungen entwickelt, unter anderem durch den Einsatz des \textit{Producer-Consumer- Dataflow-Patterns} in einem multithreaded Kontext, um Datenflüsse parallelisiert zu verarbeiten und gleichzeitig eine optimale Ressourcennutzung zu gewährleisten. Zusätzlich konnte ich durch die Unterstützung bei der Integration von Forschungsdatenmodellen auf der \href{https://coscine.rwth-aachen.de/login?redirect=/}{\textcolor{green!37!black}{Coscine}}-Plattform wertvolle Erfahrungen im Bereich Datenmanagement und Cloud-Technologien sammeln.


% Während meines Studiums der Angewandten Mathematik und Informatik sowie meiner Ausbildung zum Mathematisch-technischen Softwareentwickler habe ich fundierte Erfahrungen in der Implementierung und Optimierung von Softwarekomponenten gesammelt. So habe ich unter anderem einen \href{https://montahaee.github.io/de/projects/1_project/}{\textcolor{green!37!black}{Gitter-Generator}} für \textit{NX CAD} in C\# entwickelt und einen UML-konformen Workflow-Management für additive Fertigungsprozesse in 
%  der Automobilbaulteile \href{https://montahaee.github.io/de/projects/3_project/}{\textcolor{green!37!black}{entworfen}}, wobei ich das \textit{Adapter-} bzw. \textit{Mediator-Pattern} erfolgreich angepasst angewendet habe. Meine Fähigkeiten in der Java-Programmierung habe ich insbesondere im Rahmen der \href{https://github.com/montahaee/Informatics}{\textcolor{green!37!black}{Vorbereitung}} auf die praktische Durchführung der IHK-Abschlussprüfung zur Entwicklung eines \href{https://github.com/montahaee/SS2023/tree/main/softwareProduct}{\textcolor{green!37!black}{Softwareprodukts}} vertieft. Hierbei habe ich das \textit{Producer-Consumer-Dataflow-Pattern} in einem multithreaded Kontext eingesetzt, um die parallele Verarbeitung von Datenflüssen zu optimieren und gleichzeitig eine effiziente Ressourcennutzung zu gewährleisten.
% Während meiner MATSE-Ausbildung habe ich mich intensiv mit der Entwicklung von Softwarelösungen in C\#, Java und Python beschäftigt. 
 

% Während betrieblicher Projekte  als MATSE-Azubi konnte ich weiterhin grundlagende Kenntnisse in der Analyse von Prozessen, der Implementierung und Optimierung von Softwarekomponenten sammeln, vor allem 
% Während meiner MATSE-Ausbildung habe ich mich intensiv mit der Entwicklung von Softwarelösungen in C\#, und Python beschäftigt. 
% Insbesondere die Anwendung von \href{https://montahaee.github.io/blog/2023/dive-into-design-pattern/}{\textcolor{cyan!50!black}{Entwurfsmustern}}, die zur Strukturierung, Flexibilität und Wiederverwendbarkeit des Codes beitragen, sowie die Durchführung von Unit-Tests, die zur Sicherung der Funktionalität einzelner Softwarekomponenten und damit zu robusten und wartbaren Softwarelösungen beitragen, gehörten zu meinen Aufgaben. Darüber hinaus habe ich mein Wissen im objektiven Datenmanagement vor allem durch meine Mitarbeit an der Integration von Forschungsdatenmodellen auf der \href{https://coscine.rwth-aachen.de/login?redirect=/}{\textcolor{green!37!black}{Coscine}}-Plattform erweitert, was mir einen tiefen Einblick in agile Arbeitsmethoden und moderne Cloud-Technologien ermöglicht hat. 

% Weiterhin habe ich mich im Zuge einer Vorlesung meines AMI-Studiums, nämlich Webengineering, intensiv mit \textit{HTML, CSS} und \textit{JavaScript} beschäftigt. Später habe ich diese Kenntnisse bei der Entwicklung meiner eigenen Website mehr verfeinert. Zusätzlich konnte ich in Verbindung mit einem anderen Modul zu (relationalen)Datenbanken  nicht nur Erfahrung mit Datenbankmanagementsystemen wie \textit{MySQL} und \textit{MariaDB} sondren auch mit \href{https://github.com/montahaee/DB/tree/master/ex_am_cise/docker}{\textcolor{green!37!black}{Docker}} Container bzw. Docker-Compose aufbauen. Ferner konnte ich während der Teilnahme an den Fächern \href{https://github.com/montahaee/DS}{\textcolor{green!37!black}{Data Science}} und \href{https://montahaee.github.io/assets/pdf/en/projects/heart_disease_Prediction.pdf}{\textcolor{green!37!black}{Analyse}} meines AMI-Studiums meine Kenntnisse in Python weiter vertiefen, wobei ich mich intensiv mit statistischen und datengetriebenen Modellen sowie Methoden der explorativen Datenanalyse auseinandergesetzt habe.

% Weiterhin habe ich mich im Zuge einer Vorlesung meines AMI-Studiums, nämlich Webengineering, intensiv mit \textit{HTML, CSS} und \textit{JavaScript} beschäftigt. Später habe ich diese Kenntnisse bei der Entwicklung meiner eigenen Website weiter verfeinert. Zusätzlich konnte ich in Verbindung mit einem anderen Modul zu (relationalen) Datenbanken nicht nur Erfahrung mit Datenbankmanagementsystemen wie \textit{MySQL} und \textit{MariaDB} sammeln, sondern auch mit \href{https://github.com/montahaee/DB/tree/master/ex_am_cise/docker}{\textcolor{green!37!black}{Docker}} Containern bzw. Docker-Compose. Ferner konnte ich während der Teilnahme an den Fächern \href{https://github.com/montahaee/DS}{\textcolor{green!37!black}{Data Science}} und \href{https://montahaee.github.io/assets/pdf/en/projects/heart_disease_Prediction.pdf}{\textcolor{green!37!black}{Analyse}} meines AMI-Studiums meine Kenntnisse in Python weiter vertiefen, wobei ich mich intensiv mit statistischen und datengetriebenen Modellen sowie Methoden der explorativen Datenanalyse auseinandergesetzt habe.

% Weiterhin habe ich mich im Rahmen einer Vorlesung meines AMI-Studiums, nämlich Webengineering, intensiv mit \textit{HTML}, \textit{CSS} und \textit{JavaScript} beschäftigt. Später habe ich diese Kenntnisse bei der Entwicklung meiner eigenen Website weiter verfeinert. Zusätzlich konnte ich in Verbindung mit einem anderen Modul zu (relationalen) Datenbanken nicht nur Erfahrung mit Datenbankmanagementsystemen wie \textit{MySQL} und \textit{MariaDB} sammeln, sondern auch mit \href{https://github.com/montahaee/DB/tree/master/ex_am_cise/docker}{\textcolor{green!37!black}{Docker}}-Containern bzw. Docker-Compose. Ferner konnte ich während der Teilnahme an den Fächern \href{https://github.com/montahaee/DS}{\textcolor{green!37!black}{Data Science}} und \href{https://montahaee.github.io/assets/pdf/en/projects/heart_disease_Prediction.pdf}{\textcolor{green!37!black}{Analyse}} meines AMI-Studiums meine Kenntnisse in Python weiter vertiefen, wobei ich mich intensiv mit statistischen und datengetriebenen Modellen sowie Methoden der explorativen Datenanalyse auseinandergesetzt habe.
% Besonders schätze ich an Ihrem Programm die Möglichkeit, in verschiedenen IT-Bereichen zu rotieren und dabei sowohl in technische als auch in architektonische Themen wie Enterprise-Architecture und Webanwendungen einzutauchen, denn diese Vielfalt ermöglicht es mir, meine bisherigen Kenntnisse gezielt zu erweitern. Mein \href{https://montahaee.github.io/repositories/}{\textcolor{cyan!50!black} {Hintergrund}} in der Entwicklung von Softwarelösungen und das tiefgehende Verständnis für Cloud-Technologien bieten eine solide Basis, um mich in diesen Bereichen weiterzuentwickeln. Zudem reizt mich die Gelegenheit, technische Lösungen nicht nur zu implementieren, sondern auch auf konzeptioneller Ebene zu gestalten, was meine Begeisterung für Architekturfragen im IT-Bereich widerspiegelt. Diese Kombination aus Praxisnähe und strategischer Arbeit sehe ich als idealen Weg, um meine Fähigkeiten nachhaltig zu stärken und gleichzeitig innovative Lösungen für anspruchsvolle IT-Herausforderungen zu entwickeln.

% Während meines Studiums der Angewandten Mathematik und Informatik sowie meiner Ausbildung zum Mathematisch-technischen Softwareentwickler habe ich fundierte Kenntnisse in der Analyse von Prozessen, der Implementierung und Optimierung von Softwarekomponenten sammeln. Vor allem habe ich mich in die Entwicklung von Softwarelösungen eingearbeitet, wie zum Beispiel den \href{https://montahaee.github.io/de/projects/1_project/}{\textcolor{green!37!black}{Gitter-Generator}} für \textit{NX CAD} in C\# und \href{https://montahaee.github.io/de/projects/3_project/}{\textcolor{green!37!black}{Entwurf}} eines UML-konformen Workflow-Managements in der Automobilbaulteile. Dabei habe ich Entwurfsmuster wie das \textit{Adapter-} und \textit{Mediator-Pattern} erfolgreich angepasst. Meine Fähigkeiten in der Java-Programmierung habe ich insbesondere im Rahmen der \href{https://github.com/montahaee/Informatics}{\textcolor{green!37!black}{Vorbereitung}} auf die praktische Durchführung der IHK-Abschlussprüfung zur Entwicklung eines \href{https://github.com/montahaee/SS2023/tree/main/softwareProduct}{\textcolor{green!37!black}{Softwareprodukts}} vertieft. Hierbei habe ich das \textit{Producer-Consumer-Dataflow-Pattern} in einem multithreaded Kontext eingesetzt, um die parallele Verarbeitung von Datenflüssen zu optimieren und gleichzeitig eine effiziente Ressourcennutzung zu gewährleisten.

% Vor aller Implementierung wurden Use Case-Diagramme erstellt, um die funktionale Anforderungen zu visualisieren. In der Feindesignphase folgten Verhaltens- und Aktivitätsdiagramme, die die Abläufe im System präzisierten. Schließlich wurde die Lösung in UML-Diagrammen modelliert. Dabei trug der Einsatz von entsprechenden \href{https://montahaee.github.io/blog/2023/dive-into-design-pattern/}{\textcolor{cyan!50!black}{Entwurfsmustern}} dazu bei, den jeweiligen Entwicklungsprozess zu beschleunigen, indem bewährte Entwicklungsparadigmen genutzt wurden, die halfen, potenzielle Probleme frühzeitig aufzudecken. %und zu vermeiden.
% Nach Implementierung jeder Methode wurden Unit Tests durchgeführt, um Fehler frühzeitig zu erkennen und die Qualität sowie Zuverlässigkeit der Software sicherzustellen. 



% Vor aller Implementierung wurden Use Case-Diagramme erstellt, um die funktionale Anforderungen zu visualisieren. In der Feindesignphase folgten Verhaltens- und Aktivitätsdiagramme, die die Abläufe im System präzisierten. Schließlich wurde die Lösung in UML-Diagrammen modelliert. Dabei trug der Einsatz von entsprechenden \href{https://montahaee.github.io/blog/2023/dive-into-design-pattern/}{\textcolor{cyan!50!black}{Entwurfsmustern}} dazu bei, den jeweiligen Entwicklungsprozess zu beschleunigen, indem bewährte Entwicklungsparadigmen genutzt wurden, die halfen, potenzielle Probleme frühzeitig aufzudecken. %und zu vermeiden.
% Nach Implementierung jeder Methode wurden Unit Tests sowohl für Blackbox- als auch für Whitebox-Tests durchgeführt, um Fehler frühzeitig zu erkennen und die Qualität sowie Zuverlässigkeit der Software sicherzustellen. 

% Darüber hinaus habe ich meine Fähigkeiten in der Java-Programmierung vertieft, vor allem im Rahmen der \href{https://github.com/montahaee/Informatics}{\textcolor{green!37!black}{Vorbereitung}} auf praktischen Durchführung der IHK-Abschlussprüfung zur Entwicklung eines \href{https://github.com/montahaee/SS2023/tree/main/softwareProduct}{\textcolor{green!37!black}{Softwareprodukts}}. Hierbei habe ich \textit{Producer-Consumer- Dataflow Pattern} in einem multithreaded Kontext eingesetzt, um die parallele Verarbeitung von Datenflüssen zu optimieren und eine effiziente Ressourcennutzung zu gewährleisten. Zudem konnte ich durch meine Mitarbeit an der Integration von Forschungsdatenmodellen auf der \href{https://coscine.rwth-aachen.de/login?redirect=/}{\textcolor{green!37!black}{Coscine}}-Plattform wertvolle Einblicke in Objektdatenbank und Cloud-Technologien gewinnen.

% Darüber hinaus habe ich meine Fähigkeiten in der Java-Programmierung vertieft, insbesondere im Rahmen der \href{https://github.com/montahaee/Informatics}{\textcolor{green!37!black}{Vorbereitung}} auf die praktische Durchführung der IHK-Abschlussprüfung zur Entwicklung eines \href{https://github.com/montahaee/SS2023/tree/main/softwareProduct}{\textcolor{green!37!black}{Softwareprodukts}}. Hierbei habe ich das \textit{Producer-Consumer-Dataflow-Pattern} in einem multithreaded Kontext eingesetzt, um die parallele Verarbeitung von Datenflüssen zu optimieren und eine effiziente Ressourcennutzung zu gewährleisten. 
% Zudem konnte ich durch meine Mitarbeit an der Integration von Forschungsdatenmodellen auf der \href{https://coscine.rwth-aachen.de/login?redirect=/}{\textcolor{green!37!black}{Coscine}}-Plattform wertvolle Einblicke in Objektdatenbanken und Cloud-Technologien gewinnen. Außerdem habe ich bei allen Implementierungsaufgaben Unit Tests sowohl für Blackbox- als auch für Whitebox-Tests durchgeführt und am Ende Batch-Skripte erstellt. Diese Erfahrungen haben mein Interesse an Softwaretests und Qualitätssicherung weiter verstärkt.

% Darüber hinaus habe ich meine Fähigkeiten in der Java-Programmierung vertieft, insbesondere im Rahmen der \href{https://github.com/montahaee/Informatics}{\textcolor{green!37!black}{Vorbereitung}} auf die praktische Durchführung der IHK-Abschlussprüfung zur Entwicklung eines \href{https://github.com/montahaee/SS2023/tree/main/softwareProduct}{\textcolor{green!37!black}{Softwareprodukts}}. Hierbei habe ich das \textit{Producer-Consumer-Dataflow-Pattern} in einem multithreaded Kontext eingesetzt, um die parallele Verarbeitung von Datenflüssen zu optimieren und gleichzeitig eine effiziente Ressourcennutzung zu gewährleisten. Bei allen Implementierungsaufgaben wurden Unit Tests sowohl für Blackbox- als auch für Whitebox-Tests durchgeführt, um Fehler frühzeitig zu erkennen und die Qualität sowie Zuverlässigkeit der Software sicherzustellen. Am Ende der jeweiligen Softwareentwicklungen wurden Batch-Skripte erstellt. 


% Darüber hinaus habe ich meine Fähigkeiten in der Java-Programmierung vertieft, insbesondere im Rahmen der \href{https://github.com/montahaee/Informatics}{\textcolor{green!37!black}{Vorbereitung}} auf die praktische Durchführung der IHK-Abschlussprüfung zur Entwicklung eines \href{https://github.com/montahaee/SS2023/tree/main/softwareProduct}{\textcolor{green!37!black}{Softwareprodukts}}. Hierbei habe ich das \textit{Producer-Consumer-Dataflow-Pattern} in einem multithreaded Kontext eingesetzt, um die parallele Verarbeitung von Datenflüssen zu optimieren und gleichzeitig eine effiziente Ressourcennutzung zu gewährleisten. Bei allen Implementierungsaufgaben wurden Unit Tests sowohl für Blackbox- als auch für Whitebox-Tests durchgeführt, um Fehler frühzeitig zu erkennen und die Qualität sowie Zuverlässigkeit der Software sicherzustellen. Dabei trug der Einsatz von entsprechenden \href{https://montahaee.github.io/blog/2023/dive-into-design-pattern/}{\textcolor{cyan!50!black}{Entwurfsmustern}} dazu bei, den jeweiligen Entwicklungsprozess zu beschleunigen, indem bewährte Entwicklungsparadigmen genutzt wurden, die halfen, potenzielle Probleme frühzeitig zu identifizieren und zu vermeiden.


% Weiterhin habe ich im Zuge einer Vorlesung meines Studiums nämlich Webenginering einseits und  später Entwicklung meiner jekyll-basierten Website anderseits intesive mit HTML, CSS und JavaScript auseindergesetzt.

% Weiterhin habe ich mich im Zuge einer Vorlesung meines Studiums, nämlich Webengineering, intensiv mit \textit{HTML, CSS} und \textit{JavaScript} auseinandergesetzt. Später habe ich diese Kenntnisse bei der Entwicklung meiner Jekyll-basierten Website weiter vertieft.

% Weiterhin konnte ich durch meine Mitarbeit an der Integration von Forschungsdatenmodellen auf der \href{https://coscine.rwth-aachen.de/login?redirect=/}{\textcolor{green!37!black}{Coscine}}-Plattform nicht nur wertvolle Einblicke in Objektdatenbanken und Cloud-Technologien, sondern auch in die Zusammenarbeit in interdisziplinären Teams gewinnen. Die Arbeit in einem agilen Umfeld, das iterative und inkrementelle Prozesse kombiniert, ermöglichte es mir, flexibel auf Veränderungen zu reagieren und die Software qualitativ zu verbessern. Ferner habe ich mich im Zuge einer Vorlesung meines Studiums, nämlich Webengineering, intensiv mit \textit{HTML, CSS, JavaScript} und \textit{JSON} auseinandergesetzt. Später habe ich diese Kenntnisse bei der Entwicklung meiner eigenen Website weiter vertieft. Zusätzlich konnte ich in Verbindung mit einer anderen Vorlesung zu (relationalen)Datenbanken Erfahrung mit \href{https://github.com/montahaee/DB/tree/master/ex_am_cise/docker}{\textcolor{green!37!black}{Docker}} Container bzw. Docker-Compose vertieft. Dabei habe ich virtuelle Datenbankserver konfiguriert, wobei die Konfiguration auf YAML-Dateien basierte. 

% Weiterhin konnte ich durch meine Mitarbeit an der Integration von Forschungsdatenmodellen mit Python auf der \href{https://coscine.rwth-aachen.de/login?redirect=/}{\textcolor{green!37!black}{Coscine}}-Plattform nicht nur wertvolle Einblicke in Objektdatenbanken und Cloud-Technologien, sondern auch in die Zusammenarbeit in interdisziplinären Teams gewinnen. 
% Die Arbeit in einem agilen Umfeld, das iterative und inkrementelle Prozesse kombiniert, ermöglichte es mir, flexibel auf Veränderungen zu reagieren und die Software qualitativ zu verbessern. 
% Meine Kenntnisse in Python konnte ich während der Teilnahme an den Fächern \href{https://github.com/montahaee/DS}{\textcolor{green!37!black}{Data Science}} und \href{https://montahaee.github.io/assets/pdf/en/projects/heart_disease_Prediction.pdf}{\textcolor{green!37!black}{Analyse}} meines Studiums weiter vertiefen, wobei ich mich intensiv mit statistischen und datengetriebenen Modellen sowie Methoden der explorativen Datenanalyse auseinandergesetzt habe.


 %Dabei habe ich virtuelle Datenbankserver konfiguriert, wobei die Konfiguration auf YAML-Dateien basierte. 

% Weiterhin konnte ich durch meine Mitarbeit an der Integration von Forschungsdatenmodellen auf der \href{https://coscine.rwth-aachen.de/login?redirect=/}{\textcolor{green!37!black}{Coscine}}-Plattform nicht nur wertvolle Einblicke in Objektdatenbanken und Cloud-Technologien, sondern auch in die Zusammenarbeit in interdisziplinären Teams gewinnen. Die Arbeit in einem agilen Umfeld, das iterative und inkrementelle Prozesse kombiniert, ermöglichte es mir, flexibel auf Veränderungen zu reagieren und die Software qualitativ zu verbessern. Ferner habe ich mich im Zuge einer Vorlesung meines Studiums, nämlich "Webengineering", intensive mit \textit{HTML, CSS, JavaScript} und \textit{JSON} auseinandergesetzt und diese Kenntnisse später bei der Entwicklung meiner eigenen Website weiter vertieft. 
% Zusätzlich konnte ich in einer anderen Vorlesung zu (relationalen) Datenbanken erste Erfahrungen mit \textit{MySQL, MariaDB} und ergänzend mit \href{https://github.com/montahaee/DB/tree/master/ex_am_cise/docker}{\textcolor{green!37!black}{Docker}} bzw. Docker-Compose sammeln, indem ich virtuelle Datenbankserver auf der Basis von YAML-Dateien konfiguriert habe.

% Meine Erfahrung mit \href{https://github.com/montahaee/DB/tree/master/ex_am_cise/docker}{\textcolor{green!37!black}{Docker}} Container bzw. Docker-Compose habe ich in Verbindung mit einer Vorlesung zu (relationalen)Datenbanken vertieft. Dabei habe ich virtuelle Datenbankserver konfiguriert, wobei die Konfiguration auf YAML-Dateien basierte. 

% Meine Erfahrung mit \href{https://github.com/montahaee/DB/tree/master/ex_am_cise/docker}{\textcolor{green!37!black}{Docker}} Container bzw. Docker-Compose habe ich im Zuge einer Vorlesung zu (relationalen)Datenbanken vertieft. Dabei habe ich virtuelle Datenbankserver konfiguriert, wobei die Konfiguration auf YAML-Dateien basierte. Weiterhin konnte ich durch meine Mitarbeit an der Integration von Forschungsdatenmodellen auf der \href{https://coscine.rwth-aachen.de/login?redirect=/}{\textcolor{green!37!black}{Coscine}}-Plattform nicht nur wertvolle Einblicke in Objektdatenbanken und Cloud-Technologien, sondern auch in die Zusammenarbeit in interdisziplinären Teams gewinnen.

% Zudem konnte ich durch meine Mitarbeit an der Integration von Forschungsdatenmodellen auf der \href{https://coscine.rwth-aachen.de/login?redirect=/}{\textcolor{green!37!black}{Coscine}}-Plattform wertvolle Einblicke in Objektdatenbanken und Cloud-Technologien gewinnen. Außerdem habe ich  Diese Erfahrungen haben mein Interesse an Softwaretests und Qualitätssicherung weiter verstärkt.

% Trotz der Teilnahme und Absolvierung von Fächern wie Data Science und Analyse führten meine Neugier und mein Wille, zu einem intensiven Studium von Maschinellem Lernen, das ich sowohl in einem Seminar als auch in meiner Bachelorarbeit vertieft habe. Diese Erfahrungen ermöglichen es mir, kreative Ansätze für datengetriebene Probleme zu entwickeln und innovative Algorithmen zu konzipieren.


% Trotz der Teilnahme und Absolvierung von Fächern wie Data \href{https://github.com/montahaee/DS}{\textcolor{green!37!black}{Science}} und \href{https://montahaee.github.io/assets/pdf/en/projects/heart_disease_Prediction.pdf}{\textcolor{green!37!black}{Analyse}} führte mich meine Neugier und der Wille zur ständigen Weiterentwicklung zu einem intensiven Studium von Maschinellem Lernen, das ich in meiner  \href{https://montahaee.github.io/assets/pdf/de/projects/JLL1_final.pdf}{\textcolor{green!37!black}{Seminar-}} und \href{https://montahaee.github.io/assets/pdf/de/projects/myFH_Bachelorthesis.pdf}{\textcolor{green!37!black}{Bachelorarbeit}} vertieft habe. Diese Kenntnisse ermöglichen es mir, kreative Ansätze für datengetriebene Probleme zu entwickeln und innovative Algorithmen zu konzipieren.

% Mein Interesse gilt vorallem agilen Methoden, die iterative und inkrementelle Ansätze miteinander kombinieren und mir die Möglichkeit bieten, flexibel auf Veränderungen zu reagieren und die Software kontinuierlich zu verbessern. Dabei schätze ich auch den Einsatz von Entwurfsmustern, die zur Entwicklung nachhaltiger Softwareprodukte beitragen, indem sie sowohl die Effizienz steigern als auch die Fehleranfälligkeit reduzieren. Die Wiederverwendbarkeit und Wartbarkeit der entwickelten Komponenten ermöglicht zudem langfristige ökonomische und ökologische Vorteile.

% Meine \href{https://montahaee.github.io/de/projects/#Mathematik}{\textcolor{green!37!black}{mathematische}} Ausbildung und mein analytisches Denken erlauben es mir, komplexe versicherungsmathematische Modelle zu verstehen und zu implementieren, um die Entwicklung von Softwarelösungen zu unterstützen, die Risiken durch statistische Analysen und prädiktive Modellierung effizient bewerten, Prämien basierend auf Risikoabwägungen und historischen Daten kalkulieren und langfristige Versicherungsstrategien durch Optimierungstechniken und stochastische Simulationen verbessern.

% Zusätzlich konnte ich durch die Unterstützung bei der Integration von Forschungsdatenmodellen auf der Coscine-Plattform wertvolle Erfahrungen im Bereich Datenmanagement und Cloud-Technologien sammeln.

% Zusätzlich habe ich praktische Erfahrungen in der Implementierung und Optimierung von Softwarekomponenten gesammelt, wie etwa in der Entwicklung eines \href{https://montahaee.github.io/de/projects/1_project/}{\textcolor{violet}{Gitter-Generators}} und dem \href{https://montahaee.github.io/de/projects/3_project/}{\textcolor{violet}{Entwurf}} eines UML-konformen Workflow-Managements für additive Fertigungsprozesse. In diesen Projekten wurden das Adapter bzw. das Mediator Pattern angewandt. Diese Projekte haben mir verdeutlicht, wie wichtig Clean Code und eine durchdachte Softwarearchitektur sind, und sie haben meine Leidenschaft für die Softwareentwicklung weiter gestärkt.

% Zusätzlich habe ich praktische Erfahrungen in der Implementierung und Optimierung von Softwarekomponenten gesammelt, wie etwa bei der Entwicklung eines \href{https://montahaee.github.io/de/projects/1_project/}{\textcolor{violet}{Gitter-Generators}} und dem \href{https://montahaee.github.io/de/projects/3_project/}{\textcolor{violet}{Entwurf}} eines UML-konformen Workflow-Managements für additive Fertigungsprozesse. Dabei kamen das \textit{Adapter-} und das \textit{Mediator-Pattern} zum Einsatz. Diese Erfahrungen haben mir gezeigt, wie wichtig Clean Code und eine durchdachte Softwarearchitektur sind, und sie haben meine Leidenschaft für die Softwareentwicklung weiter gestärkt.

% Darüber hinaus habe ich mich im Modul Webengineering meines AMI-Studiums sowie durch die Erstellung meiner eigenen Website intensiv mit Webtechnologien beschäftigt. Meine Erfahrung mit \href{https://github.com/montahaee/DB/tree/master/ex_am_cise/docker}{\textcolor{violet}{Docker}} bzw. Docker-Compose habe ich im Rahmen einer Vorlesung zu Datenbanken vertieft. Dabei habe ich virtuelle Datenbankserver konfiguriert, wobei die Konfiguration auf YAML-Dateien basierte.
% Während meiner Ausbildung und meines Studiums habe ich wertvolle Erfahrungen in der Softwareentwicklung gesammelt. In Projekten wie der Entwicklung einer Gitter-Generator-Bibliothek für NX CAD in C\# sowie im \href{https://github.com/montahaee/WIL_Journal/blob/46ed320d2c351ff806c1a1b379834667a49d367b/praxisbericht-1_morteza_montahaee-Matr_Nr_3262962.pdf}{\textcolor{violet}{Entwurf}} von Workflow-Management für additive Prozesse in einer UML-konformen Ablaufmodellierung (\href{https://dap-aachen.de/2022-06-22-idam}{\textcolor{violet}{IDAM}}) habe ich mein Wissen in der agilen Softwareentwicklung vertieft und meine Fähigkeiten in der interdisziplinären Zusammenarbeit erweitert. Diese Erfahrungen haben mir gezeigt, wie wichtig eine strukturierte Herangehensweise und gute Kommunikation in der Softwareentwicklung sind, was mich zu einem passenden Kandidaten für die ausgeschriebene Position macht.

% Während meiner Ausbildung und meines Studiums habe ich mich intensiv mit der Entwicklung und Optimierung von Softwareprojekten beschäftigt, darunter die Entwicklung einer Gitter-Generator-Bibliothek für NX CAD in C\#. Außerdem habe ich an der Entwicklung von Forschungsdatenmodellen auf der Coscine-Plattform mitgewirkt, wodurch ich wertvolle Einblicke in das agile Projektmanagement und die Zusammenarbeit in interdisziplinären Teams gewinnen konnte.

% Die Arbeit mit modernen Datenformaten wie XML und JSON gehört ebenfalls zu meinen bisherigen Erfahrungen. Im Rahmen der Entwicklung von Forschungsdatenmodellen auf der Coscine-Plattform habe ich Python-Funktionen geschrieben, um XML-basierte Job-Dateien von \href{https://www.eos.info/de-de}{\textcolor{violet}{EOS-Printern}} zu analysieren, beispielsweise um Koordinatendaten jeder Schicht eines 3D-Objekts zu extrahieren und in die Datenbank zu integrieren. Darüber hinaus habe ich mich im Modul Webengineering meines AMI-Studiums sowie bei der Erstellung meiner eigenen Website intensiv mit JSON beschäftigt. 

% Besonders schätze ich an dieser Stelle die Möglichkeit, die komplexe Interaktion zwischen verschiedenen Stakeholdern und der IT-Abteilung kennenzulernen und zu gestalten. Mit meiner ausgeprägten Kommunikationsfähigkeit werde ich sicherstellen, dass alle beteiligten Parteien stets transparent und effektiv miteinander kommunizieren, um die Projektziele zu erreichen. Dabei ist mir strategisches Denken besonders wichtig, um klare Prioritäten zu setzen und Entscheidungen zu treffen, die sowohl die Geschäftsanforderungen als auch die technischen Möglichkeiten optimal berücksichtigen.

% Besonders schätze ich an Ihrem Traineeprogramm die Möglichkeit, mich intensiv in die Softwareentwicklung mit Java einzuarbeiten und praktische Erfahrungen in realitätsnahen Projekten zu sammeln. Die Arbeit mit modernen Technologien wie dem Spring Framework, REST und Docker entspricht genau meinen Interessen und Zielen. Ich bin hochmotiviert, meine analytischen Fähigkeiten, schnelle Auffassungsgabe und Leidenschaft für das Programmieren einzusetzen, um mich in Ihrem Team einzubringen und gemeinsam mit Ihnen innovative Lösungen zu entwickeln.

% Diese Erfahrungen haben mein Interesse an Softwaretests und Qualitätssicherung weiter verstärkt. Ich bin überzeugt, dass die Testautomatisierung ein entscheidender Faktor für die Entwicklung stabiler und effizienter Softwarelösungen ist. Automatisierte Tests ermöglichen es, Fehler frühzeitig zu erkennen und die Qualität sowie Zuverlässigkeit der Software sicherzustellen. Ebenso wichtig ist der Einsatz von Entwurfsmustern in der Softwareentwicklung. Entwurfsmuster beschleunigen den Entwicklungsprozess, indem sie bewährte Entwicklungsparadigmen bieten, die helfen, potenzielle Probleme frühzeitig zu identifizieren und zu vermeiden.

% Diese Kenntnisse ergänzen mein Interesse an iterativen und inkrementellen Ansätzen in der Softwareentwicklung, wie z.B. der Agilen Methode. Diese bieten mir die Möglichkeit, flexibel auf Veränderungen zu reagieren und die Software kontinuierlich und qualitativ zu verbessern. Ebenso schätze ich den Einsatz von Entwurfsmustern in der Softwareentwicklung sehr, da diese dazu beitragen, nachhaltige Softwarelösungen zu entwickeln, die sowohl ökonomisch (z.B. durch effiziente Gestaltung und Verringerung der Fehleranfälligkeit) als auch ökologisch (z.B. durch Wiederverwendbarkeit und Wartbarkeit) vorteilhaft sind.

% Ihr strukturiertes Karriereprogramm, das mich gezielt auf meine Aufgaben als Product Manager vorbereitet, sowie die umfangreichen Weiterbildungsmöglichkeiten und die offene, dynamische Unternehmenskultur, entsprechen genau meinen Vorstellungen von einem idealen Berufseinstieg.

% Besonders schätze ich an dieser Stelle die Möglichkeit, die Verteilung von Daten und Reports aus ERP- und CRM-Systemen zu entwickeln und dabei eng mit verschiedenen Fachabteilungen zusammenzuarbeiten. Auch wenn ich bisher keine direkte Erfahrung mit diesen Systemen habe, sehe ich dies als eine spannende Gelegenheit, mich in diesem Bereich weiterzubilden und mein Wissen zu erweitern. Die interdisziplinäre Zusammenarbeit und die kontinuierliche Optimierung von Softwarelösungen entsprechen meinem Verständnis von moderner Softwareentwicklung, die über die reine Programmierung hinaus auch Themen wie Architektur, Softwarequalität und Configuration Management umfasst. Dabei spielen für mich Entwurfsmuster eine wichtige Rolle, da sie eine effektive und skalierbare Gestaltung von Softwarearchitekturen ermöglichen.

% Besonders schätze ich an der ausgeschriebenen Stelle die Möglichkeit, an innovativen Mendix-Projekten mitzuwirken und mich intensiv mit Low-Code-Technologien und agilen Methoden auseinanderzusetzen. Ihre achtwöchige Jump-Start-Einführung klingt nach einem vielversprechenden Start, der mir helfen wird, mein Know-how in der Praxis zu erweitern und direkt in anspruchsvolle Projekte beim Kunden einzusteigen. Die Aussicht, sowohl technische als auch fachliche Teams bei der Entwicklung und Beratung neuer Applikationen zu unterstützen, entspricht meiner Motivation, die digitale Transformation aktiv mitzugestalten.

% Die Arbeit in einem agilen Umfeld, das iterative und inkrementelle Prozesse kombiniert, bietet mir die Möglichkeit, flexibel auf Veränderungen zu reagieren und die Software qualitativ zu verbessern. Die regelmäßigen Abstimmungsrunden und die enge Zusammenarbeit im Scrum-Team passen hervorragend zu meiner Arbeitsweise und meinem Bestreben, effiziente und innovative Lösungen zu entwickeln.

% Ihr Fokus auf lebenslanges Lernen und die umfassenden Trainings- und Weiterbildungsmöglichkeiten decken sich hervorragend mit meinen beruflichen Zielen. Dazu plane ich, dank der Möglichkeit der Anwendungsänderung im Mathematikstudium, mein abgebrochenes Masterstudium in einem moderaten Tempo und in einer beruflich relevanten Richtung abzuschließen, um meine Qualifikationen weiter zu stärken.

% Ich freue mich darauf, in Ihrem Unternehmen meine Fähigkeiten weiter auszubauen und mich aktiv in die Entwicklung neuer Microservices und der Hauptapplikation HR WORKS einzubringen. Dabei schätze ich insbesondere die Möglichkeit, mit einem erfahrenen Mentor zusammenzuarbeiten und direkte Rückmeldungen zu meiner Arbeit zu erhalten, um meine beruflichen Ziele zu erreichen.
% Ich freue mich darauf, durch Ihre umfassende Begleitung und die Unterstützung durch Mentoren meine Fähigkeiten weiter auszubauen und aktiv an der Digitalisierung sowie Nachhaltigkeit mitzuwirken.
% Ich freue mich darauf, mich in einem unbefristeten Rahmen innovative Banking-Lösungen mitzugestalten und dabei ständig neue Herausforderungen zu meistern. Dabei lege ich besonders großen Wert auf die Möglichkeit, mit einem erfahrenen Mentor zusammenzuarbeiten und direkte Rückmeldungen zu meiner Arbeit zu erhalten, um meine beruflichen Ziele zu erreichen.

% Die Aussicht, in einem dynamischen Team spannende Softwareprojekte zu realisieren und gleichzeitig die Möglichkeit zu haben, durch interne und externe Weiterbildungen mein Know-how kontinuierlich zu erweitern, finde ich besonders reizvoll. Die Werte Bechtle, wie Vertrauen, Verlässlichkeit und eine offene Kommunikation, entsprechen meinen persönlichen Vorstellungen von einer produktiven und angenehmen Arbeitsumgebung.

% Ich bin begeistert von der Anwendung neuer Technologien und Tools und setze diese leidenschaftlich in meinen Projekten ein. So habe ich beispielsweise bei der Entwicklung eines Lattice-Generators als CAD-Bibliothek für NX \textit{\href{https://en.wikipedia.org/wiki/Protocol_Buffers}{\textcolor{cyan!50!black}{Protobuf}}} zur Datenserialisierung integriert, um die Effizienz der Datentransfers zu steigern und die Leistungsfähigkeit bei der Verarbeitung komplexer Strukturen deutlich zu verbessern.


% Mein Interesse an modernen Frameworks wie Spring Boot und JPA motiviert mich besonders, mein Wissen auf diesem Gebiet zu vertiefen. Ebenso sehe ich die Themen Datenmodellierung und API-Design als spannende Herausforderungen, bei denen ich meine analytischen und konzeptionellen Fähigkeiten weiter ausbauen möchte. Die enge Zusammenarbeit mit Kollegen und der Austausch zu innovativen Entwicklungen sind für mich essenziell, um qualitativ hochwertige und zukunftssichere Softwarelösungen zu schaffen.

% Ich freue mich darauf, mich in einem innovativen und vielseitigen Unternehmen wie REMONDIS weiterzuentwickeln und die Softwarelösungen mitzugestalten, die durch ihre hohe Qualität und Anpassungsfähigkeit bereits zahlreiche Kunden begeistern. Dabei lege ich besonders großen Wert auf die Möglichkeit, mit einem erfahrenen Mentor zusammenzuarbeiten und direkte Rückmeldungen zu meiner Arbeit zu erhalten, um meine beruflichen Ziele zu erreichen.

Diese vielseitigen Erfahrungen qualifizieren mich, datenbasierte Einblicke für Ihr Team beizusteuern und gemeinsam zur Weiterentwicklung  und innovative Lösungen der Energiedatenanalyse beizutragen. Über die Möglichkeit, meine Fähigkeiten und Motivation in einem persönlichen Gespräch näher vorzustellen, würde ich mich sehr freuen.
% Ich freue mich darauf, meine bisherigen Kenntnisse in Ihrem Unternehmen einzubringen, mich weiterzuentwickeln und zu beweisen, dass ich den Anforderungen gewachsen bin. Außerdem wäre Ein \textbf{\textcolor{green!32!black}{Einstieg}} für mich frühestens ab dem \textbf{\textcolor{green!32!black}{01.10.2024}} möglich.

% https://www.pflegehilfe.org/stellenangebote/1658776
% https://www.haeger-consulting.de/job/trainee-java/